% !PRINTED 46
\weberpage{46}

If the first equation (5) is multiplied by $B_0^{(1)}$, the second by $-B_1^{(1)}$, and they are added, one obtains, by the formula system mentioned (applied to $n=1$),
\[
 -M_1=B_0^{(1)}f(0)-B_1^{(1)}f(1),
\]
and so in general, if the first $\nu$ equations (5) are multiplied in order by
\[
 B_0^{(\nu)},\quad -B_1^{(\nu)},\quad +B_2^{(\nu)},\quad \ldots,\quad \pm B_\nu^{(\nu)}
\]
and then added,
\begin{equation}
 \pm M_\nu=B_0^{(\nu)}f(0)-B_1^{(\nu)}f(1)+B_2^{(\nu)}f(2)-\cdots\pm B_\nu^{(\nu)}f(\nu),
\tag{7}
\end{equation}
whereby, according to (4), the function $F(x)$ is determined and the problem is solved. It is clear that, so long as we make no special supposition about the values $f(0),f(1),\ldots,f(n)$, every arbitrary whole rational function of $x$ can be represented in this form.
% SOURCE-ORIGINAL-SIC (mathematical notation): The scan prints uppercase F(x) here although equations (1), (4), and the surrounding argument define lowercase f(x); both lanes remain source-faithful.

\begin{center}
{\large \S.\ 10.}

\vspace{0.5em}
{\large\gesperrt{Solution of the Interpolation Problem by Differences.}}
\end{center}

The definition of the whole rational function
\begin{equation}
 B_\nu^{(x)}={x(x-1)\cdots(x-\nu+1)\over 1\cdot2\cdot3\cdots\nu}
\tag{1}
\end{equation}
gives the relation, already proved earlier for the special case of a positive integer $x$,
\begin{equation}
 B_\nu^{(x+1)}=B_\nu^{(x)}+B_{\nu-1}^{(x)},\qquad B_0^{(x+1)}=B_0^{(x)}=1.
\tag{2}
\end{equation}
From this we obtain for the coefficients $M_0,M_1,\ldots,M_n$ a mode of determination which is much more convenient for practical calculation than the application of the formulae of the preceding paragraph.

Namely, by the formulae (4) and (6) of \S~9,
\begin{equation}
 f(x)=f(0)+M_1B_1^{(x)}+M_2B_2^{(x)}+\cdots+M_nB_n^{(x)},
\tag{3}
\end{equation}
and if in this expression we replace $x$ by $x+1$ and form the difference
\begin{equation}
 \Delta_x=f(x+1)-f(x)
\tag{4}
\end{equation}
then, with regard to (2),
\begin{equation}
 \Delta_x=M_1+M_2B_1^{(x)}+M_3B_2^{(x)}+\cdots+M_nB_{n-1}^{(x)},
\tag{5}
\end{equation}
from which follows (since (5) is an equation of the same kind as (3), except that $n-1$ has taken the place of $n$):
%CONT>>
